\chapter*{Introduction}
\markboth{Introduction}{}
\addstarredchapter{Introduction}

Vérifier l'identité d'une personne à partir d'un document papier ou scanné est une tâche que beaucoup d'organismes effectuent encore à la main : une secrétaire d'université, un agent administratif ou un employé de banque relit le document, puis recopie ou compare les informations avec ce qui est déjà enregistré. Cette tâche est répétitive, lente, et sujette à l'erreur humaine, surtout quand il faut traiter de nombreux dossiers.

L'idée de ce projet est simple : automatiser cette vérification. À partir d'une photo d'une carte nationale d'identité, d'un passeport, d'un certificat de scolarité ou d'un titre de séjour, le système doit être capable de lire automatiquement les informations (nom, prénom, date de naissance, numéro de document...) grâce à la reconnaissance optique de caractères (OCR), puis de comparer ces informations à celles d'une base de référence pour signaler toute incohérence.

Ce mémoire raconte la construction de ce système, de bout en bout. Dans une première partie, je présente le contexte du projet, le problème à résoudre, et les techniques existantes sur lesquelles je me suis appuyé (OCR, traitement d'image, comparaison de chaînes de caractères). Dans une seconde partie, je détaille le système que j'ai réellement construit : son architecture, ses choix techniques, et surtout le travail mené pour en améliorer la précision.

Ce dernier point mérite d'être annoncé dès l'introduction, car c'est la partie la plus intéressante du projet : une fois l'application fonctionnelle, j'ai constaté que la vérification se trompait beaucoup plus souvent que prévu, sans savoir précisément pourquoi. Plutôt que de corriger au hasard, j'ai construit un petit banc d'essai (benchmark) avec des images de test dont je connaissais à l'avance la bonne réponse pour chaque champ. Ce banc d'essai m'a permis de chiffrer précisément la précision du système (57,1\,\% de champs corrects au départ), puis de découvrir trois bugs cachés dans le code qui faisaient chuter cette précision sur des cas pourtant courants (document légèrement tourné, photo bruitée). Après correction, la précision mesurée est passée à 79,1\,\%, soit 22 points de mieux, ce qui montre l'intérêt de mesurer plutôt que de supposer.
